\documentclass{article}
\usepackage[T1]{fontenc}
\usepackage{lmodern}
\usepackage{graphicx}
\usepackage{amsmath,amssymb}
\usepackage[a4paper,margin=2cm]{geometry}
\usepackage[absolute,overlay]{textpos}
\setlength{\TPHorizModule}{1mm}
\setlength{\TPVertModule}{1mm}
\usepackage[indonesian]{babel}
\usepackage{hyperref}
\setlength{\emergencystretch}{2em}
% unit: Halaman 10

\begin{document}

% === Halaman 10 (python/native) ===

Algoritma Rijndael

• Yang dijelaskan di sini Rijndael 128-bit (anjang kunci  128 bit)

• Tidak seperti DES yang berorientasi bit, Rijndael beroperasi dalam orientasi 

byte.

• Rijndael tidak menggunakan jaringan Feistel seperti cipher blok lainnya.

• Enciphering dilakukan dalam sejumlah putaran (iterate cipher). Setiap 

putaran mengunakan kunci internal yang berbeda (disebut round key). 

• Enciphering melibatkan operasi substitusi dan permutasi.

\end{document}
